\thispagestyle{plain}
\noindent{\sffamily\small\bfseries\color{teal} НАУЧНЫЙ ОТЧЁТ
\hfill Ti--Al \(\cdot\) 25 ноября 2026 г.\par}
\vspace{3mm}
\noindent{\sffamily\bfseries\fontsize{23}{27}\selectfont\color{ink}
Границы \(\boldsymbol{\alpha/\alpha_2}\) в Ti--Al\\
при 600 и 1273,15 К\par}
\vspace{1.5mm}\noindent{\sffamily\large\color{muted}
Прямой расчёт на MTP при нулевом давлении, C46, CALPHAD и эксперимент\par}
\vspace{2mm}\noindent{\color{teal}\rule{\linewidth}{1pt}}\par

\noindent При \textbf{600 К} прямой классический расчёт на опубликованном MTP
даёт \(\boldsymbol{\LowAlpha}\) ат.\% Al в \(\alpha\) и
\(\boldsymbol{\LowOrdered}\) ат.\% Al в \(\alpha_2\). В основной
ячейке из 1400 атомов объём и отношение осей \(c/a\) независимы,
тепловые смещения ограничены условием \(W<0{,}40\). Условные 95\%-е
интервалы по временным блокам составляют
\([\LowAlphaLo;\,\LowAlphaHi]\) и \([\LowOrderedLo;\,\LowOrderedHi]\) ат.\% Al.

При \textbf{1273,15 К} измеренные границы составляют \(19{,}5\) и
\(23{,}0\) ат.\% Al. На этой изотерме ни прямой расчёт MTP, ни
Монте-Карло с C46 не разрешили два равновесных состава. Поэтому
численные границы MTP для этой температуры остаются неопределёнными.

\begin{center}\small
\begin{tabular}{@{}lrrr@{}}
\toprule
Модель или измерение & 600 К & 1000 К & 1273,15 К\\
\midrule
Прямой MTP, свободное \(c/a\) & \(\LowAlpha/\LowOrdered\) & 14,98/22,44 & две фазы не разрешены\\
C46, классический QHA--MC    & 4,55/24,10 & 14,90/21,92 & нет касательной\\
C46, квантовый QHA--MC       & 4,59/24,09 & 14,96/21,96 & \(\dagger\)\\
CALPHAD, Liang                & 10,93/20,04 & 14,75/21,00 & 19,35/23,19\\
Эксперимент, Kainuma и соавт. & --- & --- & 19,5/23,0\\
\bottomrule
\end{tabular}
\end{center}
\noindent{\small В ячейках приведены \(x_\alpha/x_{\alpha_2}\), ат.\% Al.
\(\dagger\) В исходном архиве квантового C46 обе границы разрешены при 1200~К,
при 1225, 1250 и 1300~К касательная с двумя различными контактами не найдена.
Прочерки в строке эксперимента означают отсутствие пары при точно этой
температуре в указанной таблице.}
